
\renewcommand{\refname}{Referències} % for article class

\begin{thebibliography}{9}

\bibitem{engibaryan2023}
R. Engibaryan, \textit{The Ternary Search Tree Data Structure}, Baeldung on Computer Science, March 17, 2023. Available: \url{https://www.baeldung.com/cs/ternary-search-tree}.

\bibitem{maabar2014}
M. Maabar, \textit{Trie Data Structure}, CVR Bioinformatics, November 17, 2014. Available: \url{https://bioinformatics.cvr.ac.uk/trie-data-structure/}.

\bibitem{sharma_nodate}
A. Sharma and T. Goyal, \textit{Real World Applications of Tries} [PDF], Academia.edu. Available: \url{https://www.academia.edu/download/104436863/Research_Paper.pdf}.

\bibitem{mehta2004}
D. P. Mehta and S. Sahni, \textit{Handbook of Data Structures and Applications}, Chap. 28: Tries, Chapman and Hall/CRC, 2004. Available: \url{https://dpvipracollege.ac.in/wp-content/uploads/2023/01/Handbook_of_Data_Structures_and_Applications_Dinesh_P_Mehtawww.ebook-dl.com_.pdf}.

\bibitem{fredkin1960}
E. Fredkin, ``Trie memory'' \textit{Communications of the ACM}, vol. 3, no. 9, pp. 490--499, Sept. 1960. Available: \url{https://doi.org/10.1145/367390.367400}.


\end{thebibliography}


\section*{Apèndixs (opcionals)}

En aquest apartat podeu incloure material addicional que consideri rellevant per complementar el vostre treball, com ara taules, gràfics, fragments de codi, exemples extres, resultats ampliats o documentació tècnica. Tot i que els apèndixs no són obligatoris, poden servir per reforçar o justificar parts de l’informe principal. Tingueu en compte que no seran necessàriament considerats en la correcció si el seu contingut no aporta valor directe a la memòria principal.
